\section{Classes}

\begin{frame}{Why classes?}
  \begin{block}{}
    \begin{itemize}
      \item All objects in Python are made from a class
      \item You do not need to know about classes to use Python
      \item But class programming is powerful
    \end{itemize}
  \end{block}
  \begin{block}{}
    \begin{itemize}
      \item Class $=$ functions $+$ variables packed together
      \item A class is a logical unit in a program
      \item A large program as a combination of appropriate units
    \end{itemize}
  \end{block}
\end{frame}

\begin{frame}[fragile]{A very simple class}
  \begin{itemize}
    \item One variable: \texttt{a}
    \item One function: \texttt{dump} for printing \texttt{a}
  \end{itemize}
\begin{lstlisting}[style=pythoncompact]
class Trivial:
    def __init__(self, a):
        self.a = a

    def dump(self):
        print(self.a)
\end{lstlisting}
  Terminology: functions are called \emph{methods} and variables are
  called \emph{attributes}.
\end{frame}

\begin{frame}[fragile]{How can we use this class?}
  First, make an \emph{instance} (object) of the class:
\begin{lstlisting}[style=pythoncompact]
t = Trivial(a=4)
t.dump()
\end{lstlisting}
  Note:
  \begin{itemize}
    \item \texttt{Trivial(a=4)} means \texttt{Trivial.\_\_init\_\_(t, 4)}
    \item \texttt{self} is an argument in \texttt{\_\_init\_\_} and
          \texttt{dump}, but not used in the calls
    \item \texttt{\_\_init\_\_} is the \emph{constructor}
    \item \texttt{t.dump()} means \texttt{Trivial.dump(t)}
          (\texttt{self} is \texttt{t})
  \end{itemize}
\end{frame}

\begin{frame}{The \texttt{self} argument}
  It takes time and experience to understand \texttt{self}!
  \begin{enumerate}
    \item \texttt{self} must always be the first argument
    \item \texttt{self} is never used in calls
    \item \texttt{self} is used to access attributes and methods
          inside methods
  \end{enumerate}
  \begin{block}{}
    \texttt{self} is confusing at first, but later it greatly helps
    understanding how classes work!
  \end{block}
\end{frame}

\begin{frame}[fragile]{A class for a mathematical function}
  Function with one independent variable $t$ and parameters $v_0$, $a$:
  \[
    s(t; v_0, a) = v_0 t + \frac{1}{2} a t^2
  \]
  Class representation:
  \begin{itemize}
    \item \texttt{v0} and \texttt{a} are attributes (data)
    \item A method to evaluate $s(t)$ as a function of \texttt{t} only
  \end{itemize}
  Usage:
\begin{lstlisting}[style=pythoncompact]
s = Distance(v0=2, a=0.5)  # create instance
v = s(t=0.2)               # compute formula
\end{lstlisting}
\end{frame}

\begin{frame}[fragile]{The class code}
\begin{lstlisting}[style=pythoncompact]
class Distance:
    def __init__(self, v0, a):
        self.v0 = v0
        self.a = a

    def __call__(self, t):
        v0, a = self.v0, self.a
        return v0*t + 0.5*a*t**2

s = Distance(v0=2, a=0.5)
v = s(t=0.2)   # actually s.__call__(t=0.2)
\end{lstlisting}
\end{frame}

\begin{frame}[fragile]{General $f(x,y,z; p_1,\ldots,p_n)$}
  \begin{itemize}
    \item The $n$ parameters $p_1,\ldots,p_n$ are attributes
    \item \texttt{\_\_call\_\_(self, x, y, z)} computes $f(x,y,z)$
  \end{itemize}
\begin{lstlisting}[style=pythonTiny]
class F:
    def __init__(self, p1, p2, ...):
        self.p1 = p1
        self.p2 = p2
        ...

    def __call__(self, x, y, z):
        # return formula with x, y, z and self.p1, ...

f = F(p1=..., p2=..., ...)
print(f(1, 4, -1))
\end{lstlisting}
\end{frame}
